% ================================================================
% MLOPS & MODEL DEPLOYMENT
% Comprehensive presentation on production ML systems
%
% Topics Covered:
% 1. MLOps Fundamentals and Lifecycle
% 2. Model Training Pipeline and Experiment Tracking
% 3. Model Deployment Strategies
% 4. Monitoring and Maintenance
% 5. CI/CD for ML
% 6. Best Practices and Case Studies
%
% Total: ~45 slides
% Author: Diogo Ribeiro
% Date: January 2025
% ================================================================

\section{MLOps \& Model Deployment}

% ================================================================
% PART 1: MLOPS FUNDAMENTALS
% ================================================================

\begin{frame}{From Notebook to Production}
\begin{table}
\small
\begin{tabular}{ll}
\toprule
\textbf{Development focus} & \textbf{Production concern} \\
\midrule
Static datasets & Live, changing data \\
Manual runs & Automated, repeatable pipelines \\
One local model & Versioned model/data/config artifacts \\
Offline evaluation & Online service and business behavior \\
Single-machine assumptions & Capacity, latency, and failure handling \\
Limited monitoring & Observability and incident response \\
\bottomrule
\end{tabular}
\end{table}

\begin{alertblock}{Production gap}
A model that performs well in development still needs reliable data contracts, serving behavior, monitoring, rollback, and operational ownership.
\end{alertblock}
\end{frame}

\begin{frame}{MLOps as an Operating Discipline}
\begin{columns}
\begin{column}{0.48\textwidth}
\textbf{Core concerns}
\begin{itemize}
  \item automation and repeatability;
  \item reproducibility across code, data, configuration, and environment;
  \item reliability and graceful failure;
  \item capacity and latency;
  \item traceability and governance;
  \item collaboration across data, ML, and platform teams.
\end{itemize}
\end{column}

\begin{column}{0.48\textwidth}
\begin{figure}
\centering
\begin{tikzpicture}[scale=0.85]
\node[draw, circle, fill=forest!20, minimum size=1.8cm] (ml) at (0,1.5) {ML};
\node[draw, circle, fill=navyblue!20, minimum size=1.8cm] (dev) at (2.3,1.5) {DevOps};
\node[draw, circle, fill=crimson!20, minimum size=1.8cm] (data) at (1.15,0) {Data Eng};
\node[draw, circle, fill=purple!40, minimum size=1.0cm] (center) at (1.15,0.9) {MLOps};
\end{tikzpicture}
\end{figure}
\end{column}
\end{columns}

\begin{block}{Not just tooling}
MLOps is the set of engineering practices and organizational controls used to operate model-based systems over time.
\end{block}
\end{frame}

\begin{frame}{The ML Lifecycle: Build}
\textbf{Model artifacts depend on more than code:}
\[
\text{Code} + \text{Data} + \text{Configuration} + \text{Environment}
\longrightarrow \text{Model artifact}.
\]

\begin{enumerate}
  \item \textbf{Data collection}: acquire and validate source data.
  \item \textbf{Data preparation}: clean, transform, and construct features.
  \item \textbf{Training}: fit candidate models under tracked configurations.
  \item \textbf{Evaluation}: compare against baselines and release criteria.
\end{enumerate}

\begin{alertblock}{Reproducibility unit}
To reproduce a model, record the source code, data/version identifiers, configuration, environment, and randomization state—not only the serialized model file.
\end{alertblock}
\end{frame}

\begin{frame}{The ML Lifecycle: Operate}
\begin{columns}
\begin{column}{0.46\textwidth}
\begin{enumerate}
  \setcounter{enumi}{4}
  \item \textbf{Deploy}: expose predictions or decisions through the required interface.
  \item \textbf{Monitor}: observe system health, inputs, outputs, labels, and business outcomes.
  \item \textbf{Update}: retrain, recalibrate, roll back, or retire when evidence justifies a change.
\end{enumerate}
\end{column}

\begin{column}{0.5\textwidth}
\begin{figure}
\centering
\begin{tikzpicture}[scale=0.78]
\node[draw, circle, fill=forest!20] (data) at (0,2) {Data};
\node[draw, circle, fill=navyblue!20] (train) at (2.2,2) {Train};
\node[draw, circle, fill=crimson!20] (deploy) at (4.4,2) {Deploy};
\node[draw, circle, fill=purple!20] (monitor) at (4.4,0) {Monitor};
\node[draw, circle, fill=gold!20] (update) at (0,0) {Update};
\draw[->, thick] (data) -- (train);
\draw[->, thick] (train) -- (deploy);
\draw[->, thick] (deploy) -- (monitor);
\draw[->, thick] (monitor) -- (update);
\draw[->, thick] (update) -- (data);
\end{tikzpicture}
\end{figure}
\end{column}
\end{columns}

\begin{block}{Feedback loop}
New data, changing requirements, delayed labels, incidents, and measured performance can all trigger review. Continuous retraining is not automatically the correct response.
\end{block}
\end{frame}

% ================================================================
% PART 2: EXPERIMENT TRACKING AND VERSIONING
% ================================================================

\begin{frame}[fragile,allowframebreaks]{Experiment Tracking}
\textbf{Managing the chaos of ML experiments}

\begin{columns}
\begin{column}{0.5\textwidth}
\textbf{The Problem:}

\begin{itemize}
\item Tried 50 experiments
\item Which hyperparameters worked best?
\item Can't reproduce last week's results
\item Lost track of what changed
\end{itemize}

\vspace{0.3cm}

\textbf{What to Track:}

\begin{enumerate}
\item \textcolor{forest}{\textbf{Code:}}
   \begin{itemize}
   \item Git commit hash
   \item Branch name
   \item Dependencies (requirements.txt)
   \end{itemize}

\item \textcolor{navyblue}{\textbf{Data:}}
   \begin{itemize}
   \item Dataset version
   \item Data hash/signature
   \item Train/val/test splits
   \end{itemize}

\item \textcolor{crimson}{\textbf{Model:}}
   \begin{itemize}
   \item Architecture
   \item Hyperparameters
   \item Random seeds
   \end{itemize}

\item \textcolor{purple}{\textbf{Metrics:}}
   \begin{itemize}
   \item Accuracy, F1, AUC
   \item Training time
   \item Model size
   \end{itemize}
\end{enumerate}
\end{column}

\begin{column}{0.5\textwidth}
\textbf{MLflow Example:}

\begin{verbatim}
import mlflow

mlflow.set_experiment("churn_prediction")

with mlflow.start_run():
    # Log parameters
    mlflow.log_param("n_estimators", 100)
    mlflow.log_param("max_depth", 10)

    # Train model
    model = RandomForestClassifier(
        n_estimators=100,
        max_depth=10
    )
    model.fit(X_train, y_train)

    # Log metrics
    accuracy = model.score(X_test, y_test)
    mlflow.log_metric("accuracy", accuracy)

    # Log model
    mlflow.sklearn.log_model(model, "model")

    # Log artifacts
    mlflow.log_artifact("confusion_matrix.png")
\end{verbatim}

\vspace{0.2cm}

\textbf{Popular Tools:}
\begin{itemize}
\item MLflow (open source)
\item Weights \& Biases
\item Neptune.ai
\item Comet.ml
\end{itemize}
\end{column}
\end{columns}
\end{frame}

\begin{frame}[fragile,allowframebreaks]{Model and Data Versioning}
\textbf{Git for machine learning assets}

\begin{columns}
\begin{column}{0.5\textwidth}
\textbf{Why Version Everything?}

\begin{itemize}
\item \textbf{Reproducibility}: Recreate any experiment
\item \textbf{Debugging}: What changed between v1 and v2?
\item \textbf{Rollback}: Revert to previous version
\item \textbf{Compliance}: Audit trail for regulations
\item \textbf{Collaboration}: Team knows what's happening
\end{itemize}

\vspace{0.3cm}

\textbf{Data Version Control (DVC):}

Like Git, but for data and models:

\begin{verbatim}
# Initialize DVC
dvc init

# Track data file
dvc add data/raw/dataset.csv

# Commit to git (only hash stored)
git add data/raw/dataset.csv.dvc
git commit -m "Add raw dataset v1"

# Push data to remote storage
dvc push

# Pull data (like git pull)
dvc pull
\end{verbatim}
\end{column}

\begin{column}{0.5\textwidth}
\textbf{DVC Benefits:}

\begin{itemize}
\item Data stored in cloud (S3, GCS, Azure)
\item Git tracks only metadata
\item Lightweight repositories
\item Easy collaboration
\item Pipeline versioning
\end{itemize}

\vspace{0.2cm}

\textbf{ML Pipeline with DVC:}

\begin{verbatim}
# dvc.yaml
stages:
  preprocess:
    cmd: python preprocess.py
    deps:
      - data/raw
    outs:
      - data/processed

  train:
    cmd: python train.py
    deps:
      - data/processed
      - train.py
    params:
      - model.n_estimators
    outs:
      - models/model.pkl
    metrics:
      - metrics.json
\end{verbatim}

\vspace{0.2cm}

\textbf{Run pipeline:}
\begin{verbatim}
dvc repro
\end{verbatim}

DVC automatically tracks dependencies and reruns only what changed!
\end{column}
\end{columns}
\end{frame}

% ================================================================
% PART 3: MODEL DEPLOYMENT
% ================================================================

\begin{frame}[fragile,allowframebreaks]{Deployment Patterns}
\textbf{How to serve model predictions}

\begin{columns}
\begin{column}{0.5\textwidth}
\textbf{1. Batch Prediction:}

\begin{itemize}
\item Process data in batches (hourly, daily)
\item Generate predictions offline
\item Store in database
\item Serve pre-computed predictions
\end{itemize}

\textbf{Use when:}
\begin{itemize}
\item[$\checkmark$] Don't need real-time
\item[$\checkmark$] Can predict in advance
\item[$\checkmark$] Large-scale scoring
\end{itemize}

\textbf{Examples:}
\begin{itemize}
\item Daily recommendation emails
\item Credit score calculations
\item Churn probability (monthly)
\end{itemize}

\vspace{0.3cm}

\textbf{2. Real-Time API:}

\begin{itemize}
\item Model as REST API
\item Synchronous: request $\to$ prediction
\item Sub-second latency required
\end{itemize}

\textbf{Use when:}
\begin{itemize}
\item[$\checkmark$] Immediate response needed
\item[$\checkmark$] User-facing applications
\item[$\checkmark$] Input varies per request
\end{itemize}
\end{column}

\begin{column}{0.5\textwidth}
\textbf{Flask API Example:}

\begin{verbatim}
from flask import Flask, request, jsonify
import joblib

app = Flask(__name__)
model = joblib.load('model.pkl')

@app.route('/predict', methods=['POST'])
def predict():
    data = request.get_json()
    features = data['features']

    prediction = model.predict([features])
    probability = model.predict_proba([features])

    return jsonify({
        'prediction': int(prediction[0]),
        'probability': float(probability[0][1])
    })

if __name__ == '__main__':
    app.run(host='0.0.0.0', port=5000)
\end{verbatim}

\vspace{0.2cm}

\textbf{3. Streaming:}

\begin{itemize}
\item Process events as they arrive
\item Kafka, Kinesis
\item Low latency, high throughput
\end{itemize}

\vspace{0.2cm}

\textbf{4. Edge Deployment:}

\begin{itemize}
\item Model on device (phone, IoT)
\item No internet required
\item Ultra-low latency
\item Privacy-preserving
\end{itemize}
\end{column}
\end{columns}
\end{frame}

\begin{frame}[fragile,allowframebreaks]{Containerization with Docker}
\textbf{Package model + dependencies + code together}

\begin{columns}
\begin{column}{0.5\textwidth}
\textbf{Why Docker?}

\begin{itemize}
\item \textbf{"Works on my machine"} $\to$ Works everywhere!
\item Reproducible environments
\item Easy deployment
\item Isolated dependencies
\item Scalable (Kubernetes)
\end{itemize}

\vspace{0.3cm}

\textbf{Dockerfile Example:}

\begin{verbatim}
FROM python:3.9-slim

# Set working directory
WORKDIR /app

# Copy requirements
COPY requirements.txt .

# Install dependencies
RUN pip install --no-cache-dir -r
    requirements.txt

# Copy application
COPY . .

# Expose port
EXPOSE 5000

# Run app
CMD ["python", "app.py"]
\end{verbatim}
\end{column}

\begin{column}{0.5\textwidth}
\textbf{Build and Run:}

\begin{verbatim}
# Build image
docker build -t my-ml-model:v1 .

# Run container
docker run -p 5000:5000 my-ml-model:v1

# Test
curl -X POST http://localhost:5000/predict \
  -H "Content-Type: application/json" \
  -d '{"features": [1.0, 2.0, 3.0]}'
\end{verbatim}

\vspace{0.2cm}

\textbf{Docker Compose (Multi-Container):}

\begin{verbatim}
# docker-compose.yml
version: '3'
services:
  model:
    build: .
    ports:
      - "5000:5000"

  redis:
    image: redis:alpine

  monitoring:
    image: prometheus
    ports:
      - "9090:9090"
\end{verbatim}

\vspace{0.2cm}

\begin{verbatim}
docker-compose up
\end{verbatim}
\end{column}
\end{columns}

\vspace{0.2cm}

\begin{alertblock}{Best Practice}
Always containerize ML models for consistent deployment across environments!
\end{alertblock}
\end{frame}

% ================================================================
% PART 4: MONITORING
% ================================================================

\begin{frame}{Production Monitoring: Model and Data Signals}
\begin{itemize}
  \item \textbf{Data quality}: schema, missingness, ranges, categorical levels, freshness.
  \item \textbf{Input distribution}: compare production data with relevant reference windows.
  \item \textbf{Prediction distribution}: watch for changes in score/rate distributions.
  \item \textbf{Observed performance}: evaluate loss, calibration, discrimination, or task-specific metrics once labels arrive.
\end{itemize}

\begin{alertblock}{Drift is evidence, not diagnosis}
A distribution shift can be harmless, and performance can fail without large marginal feature drift. Monitoring should connect signals to downstream model and business outcomes.
\end{alertblock}
\end{frame}

\begin{frame}{Production Monitoring: System and Decision Signals}
\begin{columns}
\begin{column}{0.48\textwidth}
\textbf{System health}
\begin{itemize}
  \item latency percentiles;
  \item throughput;
  \item error/timeout rates;
  \item resource saturation.
\end{itemize}
\end{column}

\begin{column}{0.48\textwidth}
\textbf{Decision/business health}
\begin{itemize}
  \item intervention or approval rates;
  \item downstream outcomes;
  \item subgroup behavior;
  \item operational costs and incidents.
\end{itemize}
\end{column}
\end{columns}

\begin{block}{Alert thresholds}
Thresholds should be calibrated to service objectives, historical variability, traffic volume, label delay, and consequence severity. There is no universal KL, accuracy-drop, latency, or error-rate cutoff.
\end{block}
\end{frame}

\begin{frame}[fragile,allowframebreaks]{CI/CD for Machine Learning}
\textbf{Continuous Integration and Deployment}

\begin{columns}
\begin{column}{0.5\textwidth}
\textbf{Traditional CI/CD:}

\begin{enumerate}
\item Code commit $\to$ Git
\item Run tests
\item Build artifact
\item Deploy to production
\end{enumerate}

\vspace{0.2cm}

\textbf{ML CI/CD Additions:}

\begin{itemize}
\item \textbf{CT:} Continuous Training
\item \textbf{CD:} Continuous Deployment
\item \textbf{CM:} Continuous Monitoring
\end{itemize}

\vspace{0.3cm}

\textbf{GitHub Actions Example:}

\begin{verbatim}
# .github/workflows/ml-pipeline.yml
name: ML Pipeline

on: [push]

jobs:
  test:
    runs-on: ubuntu-latest
    steps:
      - uses: actions/checkout@v2
      - name: Run tests
        run: pytest tests/

  train:
    needs: test
    runs-on: ubuntu-latest
    steps:
      - name: Train model
        run: python train.py

      - name: Evaluate
        run: python evaluate.py

      - name: Check metrics
        run: python check_metrics.py
\end{verbatim}
\end{column}

\begin{column}{0.5\textwidth}
\begin{verbatim}
  deploy:
    needs: train
    if: github.ref == 'refs/heads/main'
    runs-on: ubuntu-latest
    steps:
      - name: Build Docker image
        run: docker build -t model:${{
           github.sha }} .

      - name: Push to registry
        run: docker push model:${{
           github.sha }}

      - name: Deploy to staging
        run: kubectl apply -f k8s/staging/

      - name: Run integration tests
        run: pytest tests/integration/

      - name: Deploy to production
        if: success()
        run: kubectl apply -f k8s/prod/
\end{verbatim}

\vspace{0.2cm}

\textbf{Key Tests for ML:}

\begin{itemize}
\item \textbf{Data validation:} Schema, ranges
\item \textbf{Model tests:} Accuracy threshold
\item \textbf{Integration tests:} API endpoints
\item \textbf{Performance tests:} Latency SLAs
\item \textbf{Bias tests:} Fairness metrics
\end{itemize}
\end{column}
\end{columns}
\end{frame}

\begin{frame}[fragile,allowframebreaks]{Knowledge Check: MLOps}
\textbf{Test your understanding}

\vspace{0.2cm}

\textbf{Question 1:} Your model was 92\% accurate in development but 78\% in production. What's likely the issue?

\begin{enumerate}[A)]
\item The model is overfitted
\item Wrong metrics used
\item \textcolor{forest}{\textbf{Data drift - production data differs from training data}}
\item Need more training data
\end{enumerate}

\pause

\begin{block}{Explanation}
\textcolor{forest}{\textbf{Best answer: investigate distribution shift, data quality, label/metric changes, and training-serving skew.}} A production performance drop is evidence of a changed evaluation problem, not a diagnosis by itself. Input drift may be present without harming performance, while concept/label changes can matter even when marginal feature distributions look stable.
\end{block}

\pause

\vspace{0.2cm}

\textbf{Question 2:} For real-time fraud detection (need <100ms latency), which deployment pattern?

\begin{enumerate}[A)]
\item Batch prediction (daily)
\item \textcolor{forest}{\textbf{Real-time API with caching and model optimization}}
\item Edge deployment
\item Streaming with 1-hour windows
\end{enumerate}

\pause

\begin{block}{Explanation}
\textcolor{forest}{\textbf{Best answer: choose an online serving path that can meet the decision deadline.}} The model family is not determined by the latency target alone; preprocessing cost, feature freshness, hardware, batching, caching, and system architecture all contribute to end-to-end latency.
\end{block}
\end{frame}

\begin{frame}{MLOps Operating Checklist}
\begin{itemize}
  \item version code, data identifiers, model artifacts, configuration, and environment;
  \item automate reproducible training/evaluation/deployment pipelines;
  \item test data contracts, code, integrations, and model release criteria;
  \item monitor system, data, model, and business signals;
  \item maintain rollback, recovery, and incident-response procedures;
  \item document ownership, approvals, and model lifecycle decisions.
\end{itemize}
\end{frame}

\begin{frame}{MLOps: Tools Are Replaceable}
\begin{itemize}
  \item Experiment tracking, orchestration, serving, monitoring, and feature management each have many viable implementations.
  \item Architecture should depend on scale, latency, team structure, regulatory constraints, cloud/platform choices, and maintenance burden.
  \item Avoid building a large platform before the use case needs it; operational complexity is itself a cost and failure mode.
\end{itemize}

\begin{alertblock}{The durable principle}
Choose tools to support reproducibility, observability, safe release, and recovery. The product names can change without changing the engineering contract.
\end{alertblock}
\end{frame}

% ================================================================
% END OF MLOPS PRESENTATION
% ================================================================


